\scanpage{121}
aux conditions supplémentaires techniques énoncées (être une ``quasi'' résolution de $S_k$, et être ``quasi-flasque'' $\ldots$)

On note que les hom. dans \deleted{\unsure{cette}} la catégorie des objets \inserted{$C_*^{**}$} envisagés sont déjà des ``quasi-iso'' (pour $\tau_\infty$) - il y a lieu de passer à la cat. des fractions en \deleted{les} rendant inversibles tous ces hom., i.e. le \uline{groupoïde engendré} par la catégorie \uline{envisagée}. On aimerait donc prouver que ce groupoïde est connexe et simplement connexe ?
[\uline{NB} rappelons \unsure{que c'est en un} \uline{langage pr} \inserted{\unsure{en dehors}} $S_k=k\{T\}$ \unsure{bien entendu plus qu'il y ait calcul de fractions à gauche}]

Cela signifie donc

\uline{Conj.?} a) Si $C_*^{**}$ et $C'_*{}^{**}$ sont des théories admissibles, $\exists$ théorie admissible $C''_*{}^{**}$ et hom
\[
\diagram{\begin{tikzcd}
&C_*^{**}\\
C''_*{}^{**}\arrow[ur]\arrow[dr,no head]\\
&C'_*{}^{**}
\end{tikzcd}}
\]
\margin{\figure{\begin{tikzpicture}[x=1cm,y=1cm]\draw (0,0)--(0,4);\node at (-0.6,2.9) {\Huge ?};\end{tikzpicture}}}
b) Si on a deux hom
\[
\diagram{\begin{tikzcd}
C'_*{}^{**}\arrow[r,shift left=0.7ex]\arrow[r,shift right=0.7ex,no head]&C_*^{**}
\end{tikzcd}}
\]
$\exists$ hom
\[
\diagram{\begin{tikzcd}
C''_*{}^{**}\arrow[r,no head]&C'_*{}^{**}
\end{tikzcd}}
\]
qui ``égalise''.
\margin{\unsure{il y aurait} plus de chances avec calcul de fractions à droite ???}
\margin{\figure{\begin{tikzpicture}[x=1cm,y=1cm]\draw (0,0)--(1.1,3.6);\end{tikzpicture}}}
\margin{120}
